% Minimal retained preamble for lane-p2 cards taken from
% sources/logic-fl02-conservation-of-ramsey-s-theorem-f/pinned-original-source.tex
% (Conservation of Ramsey's theorem for pairs, arXiv:2402.11616v3).
%
% Copied verbatim from the author's own file:
%   line 2      : \documentclass{article}
%   lines 16-25 : the author's theorem declarations
%   lines 46-88 : the author's macro definitions for the reverse-mathematics principles
%   line 91     : \card
%   lines 104-107: the natural sums and products, which rely on mathabx (source line 103)
%   line 116    : \qt
%
% Omitted author preamble lines and why (nothing was renamed or redefined):
%   line 3 (graphicx), 7 (url), 9 (cancel): no figure, url or cancel use in the retained spans;
%   line 6 (cleveref): the assembler adds its own cleveref machinery, a second load clashes;
%   line 8 (\usepackage[utf8x]{inputenc}): inputenc with the utf8x option fights the forced
%       xelatex engine;
%   lines 28-37 (the rep@theorem machinery and \newreptheorem{maintheorem}): the retained spans
%       use no rep-theorem environment;
%   lines 40-43 (xcolor and the three author comment macros): not used by the retained spans;
%   lines 93-99 (the LS2 font encoding block for \smashtimes based on stix): the host TeX Live has
%       no stix directory and the author himself never uses the result -- line 105 maps \ntimes to
%       \dottimes, and line 110 shows \smashtimes already commented out in the source.
% mathabx is retained because \dotplus, \dottimes and \dotdiv come from it.
\documentclass{article}
\usepackage{amsmath,amssymb,amsthm}
\usepackage{mathabx}
\usepackage{cancel}
\usepackage{cleveref}

\newtheorem{theorem}{Theorem}
\numberwithin{theorem}{section}
\newtheorem{maintheorem}[theorem]{Main Theorem}
\newtheorem{lemma}[theorem]{Lemma}
\newtheorem{question}[theorem]{Question}
\newtheorem{proposition}[theorem]{Proposition}
\newtheorem{statement}[theorem]{Statement}
\newtheorem{remark}[theorem]{Remark}
\newtheorem{definition}[theorem]{Definition}
\newtheorem{corollary}[theorem]{Corollary}

\newcommand{\RCA}{\mathsf{RCA}}
\newcommand{\WKL}{\mathsf{WKL}}
\newcommand{\SRT}{\mathsf{SRT}}
\newcommand{\EM}{\mathsf{EM}}
\newcommand{\SEM}{\mathsf{SEM}}
\newcommand{\ADS}{\mathsf{ADS}}
\newcommand{\SCAC}{\mathsf{SCAC}}
\newcommand{\SADS}{\mathsf{SADS}}
\newcommand{\COH}{\mathsf{COH}}
\newcommand{\ACA}{\mathsf{ACA}}
\newcommand{\qvdash}{\operatorname{{?}{\vdash}}}
\newcommand{\nqvdash}{\operatorname{{?}{\nvdash}}}
\newcommand{\BSig}{\mathsf{B}\Sigma^0}
\newcommand{\ISig}{\mathsf{I}\Sigma^0}
\newcommand{\LSig}{\mathsf{L}\Sigma^0}
\newcommand{\BPi}{\mathsf{B}\Pi^0}
\newcommand{\IPi}{\mathsf{I}\Pi^0}
\newcommand{\LPi}{\mathsf{L}\Pi^0}
\newcommand{\IDel}{\mathsf{I}\Delta^0}
\newcommand{\BCDe}{\mathsf{BC}\Delta^0}
\newcommand{\WF}{\mathsf{WF}}
\newcommand{\RT}{\mathsf{RT}}
\newcommand{\BME}{\mathsf{BME}}
\newcommand{\TT}{\mathsf{TT}}
\renewcommand{\O}{\mathcal{O}}
\renewcommand{\L}{\mathcal{L}}
\newcommand{\C}{\mathcal{C}}
\newcommand{\R}{\mathcal{R}}
\renewcommand{\S}{\mathcal{S}}
\newcommand{\T}{\mathcal{T}}
\renewcommand{\P}{\mathcal{P}}
\newcommand{\M}{\mathcal{M}}
\newcommand{\Nc}{\mathcal{N}}
\renewcommand{\P}{\mathcal{P}}
\newcommand{\Q}{\mathcal{Q}}
\newcommand{\NN}{\mathbb{N}}
\newcommand{\PP}{\mathbb{P}}
\newcommand{\uh}[0]{{\upharpoonright}}
\newcommand{\card}{\operatorname{card}}
\newcommand{\nplus}{\dotplus}
\newcommand{\ntimes}{\dottimes}
\newcommand{\nminus}{\dotdiv}
\newcommand{\nsum}{\dot{\sum}}
\def\qt#1{``#1''}
\title{Every SEM condition decides an initial segment of the jump by an extension}
\author{Quentin Le Houérou, Ludovic Levy Patey, Keita Yokoyama}
\date{arxiv:2402.11616v3; CC BY 4.0}
\begin{document}
\maketitle
\noindent\textit{Editorial note.} The counted unit is the article's own Lemma lem:sem-wf-forcing-jump at source lines 918-920 together with the proof environment that begins on the next line (921-929); the pairing is the ordinary adjacent one.\par\medskip
\noindent\textit{Editorial note.} The retained context at source lines 853-861 is the author's own definition of trees and universal trees; the retained context at source lines 863-865 and 893-897 is the statement of the two lemmas that the counted proof invokes, namely the existence of a maximal enumeration and the finite-tree forcing lemma, the latter counted by a sibling card of this lane. They are retained as prerequisites so that the author Cref calls resolve inside the retained material and no navigation stub is needed.\par\medskip
\noindent\textit{Editorial note.} The retained preamble reproduces verbatim the author's document class at source line 2, his theorem declarations at source lines 16-25, his macro definitions at source lines 46-88 and 91, and his natural-operation macros at source lines 104-107, and it retains mathabx and cleveref; the stix-dependent LS2 block and the inputenc load with the utf8x option are not retained, the former because the host lacks stix and the latter because it fights xelatex. Nothing was renamed or redefined.\par\medskip
\noindent\textit{Editorial note.} No retained line contains a citation, so the card carries no bibliography.\par\medskip
\noindent\textit{Editorial note.} The article's conservation theorems, its COH and WKL lemmas and the other lemmas of the SEM section are not counted here and are not used as premises.\par\medskip
\noindent\textit{Allowed standard prerequisites (editorial).}\begin{itemize}
\item Monotone enumerations of maximal universal trees. For every Sigma\textasciicircum{}0\_1 formula and every condition there is a monotone enumeration of a maximal universal tree satisfying the formula, and the tree is a set of nodes closed downward under prefixes with the universality property of the SEM section; the author's own statement is retained at source lines 863-865 and the definitions of trees at 853-861.. The proof builds the enumeration for the specific formula that codes an initial segment of the jump, so both the existence statement and the notion of tree must be available.
\item A finite maximal universal tree forces a last true instance. If the enumeration is M-finite, then there is an extension of the condition and an index q such that the extension forces the formula at the index below q and forces its negation at q; retained with this card at source lines 893-897.. The proof applies this lemma to the enumeration once the enumeration is known to be M-finite, which is what turns the tree statement into the forcing statement about the jump.
\item The theory BME\_* and the bound k on the enumerated formula. Since the ordinal is positive the model satisfies WF(omega\textasciicircum{}omega) and therefore the principle BME\_* of bounded monotone enumerations, which is what makes the enumeration M-finite; moreover the specific formula is built so that it never holds at the index k+1, which bounds the enumeration.. These two facts supply the M-finiteness hypothesis and the finiteness of the tree that the proof needs before it can invoke the finite-tree lemma.
\end{itemize}\medskip
\begin{definition}
Fix a condition $(\sigma, a)$.
\begin{enumerate}
    \item[(1)] A \emph{$(\sigma, a)$-tree} is a set of $(\sigma, a)$-nodes closed downard under the prefix relation.
    \item[(2)] A $(\sigma, a)$-tree $E$ is \emph{universal} if for every~$(\sigma,a)$-node $\nu$
    which is not a leaf, there is some~$r \in M$ such that for every $M$-finite 2-partition $A_0 \sqcup A_1 = \{0, \dots, r-1\}$, there is a side~$i < 2$ and an $M$-finite $R$-transitive $\rho \subseteq A_i \cap X_a$ such that $\nu \cdot \langle i, \rho \rangle \in E$.
    \item[(3)] A $(\sigma, a)$-tree $E$ satisfies a $\Sigma^0_1(\M)$ formula if all its nodes satisfies it.
\end{enumerate}
\end{definition}

\begin{lemma}\label[lemma]{lem:sem-wf-exists-maximal-enumeration}
For every~$\Sigma^0_1(\M)$ formula $\varphi(G, x)$ and every condition $(\sigma, a)$, there is a monotone enumeration of a maximal (for inclusion) universal $(\sigma, a)$-tree~$E$ satisfying $\varphi(G, x)$.
\end{lemma} 

\begin{lemma}\label[lemma]{lem:sem-wf-finite-tree-forcing}
Let $(\sigma,a)$ be a condition, $\varphi(G, x)$ be a $\Sigma^0_1(\M)$ formula, and $E$ be an $M$-finite maximal universal $(\sigma, a)$-tree satisfying $\varphi(G, x)$.
Then there is an extension $(\tau, c) \leq (\sigma, a)$ and some~$q \in M$ such that
$(\tau, c) \Vdash \varphi(G, q-1)$ and $(\tau, c) \Vdash \neg \varphi(G, q)$.
\end{lemma}

\begin{lemma}\label[lemma]{lem:sem-wf-forcing-jump}
Let $(\sigma,a)$ be a condition. For every $k \in M$, there exists some $M$-finite $\sigma' \in 2^k$ and an extension $(\tau,b) \leq (\sigma,a)$ such that $(\tau,b) \Vdash (G \oplus Y)' \uh k = \sigma'$    
\end{lemma}

\begin{proof}
Let $\varphi(G, x)$ be the $\Sigma^0_1(Y)$ formula which holds if there exists some~$e_0 < \dots < e_{x-1} < k$ such that for every~$s < x$, $\Phi^{G \oplus Y}_{e_s}(e_s)\downarrow$.

By \Cref{lem:sem-wf-exists-maximal-enumeration}, there exists a monotone enumeration of a maximal universal $(\sigma, a)$-tree~$E$ satisfying $\varphi(G, x)$.
Moreover, since~$\varphi(G, k+1)$ never holds, the enumeration $E$ is $k$-bounded.
Since $\alpha > 0$, $\M \models \WF(\omega^\omega)$, hence $\M \models \BME_*$, so $E$ is $M$-finite.
By \Cref{lem:sem-wf-finite-tree-forcing}, there is an extension $(\tau, c) \leq (\sigma, a)$ and some~$q \in M$ such that $(\tau, c) \Vdash \varphi(G, q-1)$ and $(\tau, c) \Vdash \neg \varphi(G, q)$.
Let~$\sigma' = \{ e < k : \Phi^{\tau \oplus Y}_e(e)\downarrow \}$. Then $(\tau,b) \Vdash (G \oplus Y)' \uh k = \sigma'$.
\end{proof}

\end{document}
